% Support de présentation de la soutenance.
%
% Même moteur que le rapport — pdflatex par latexmk — et pour le même motif : la
% classe beamer et les polices Latin Modern sont présentes dans toute distribution
% complète, sans dépendre d'une police du système.
%
% Les marqueurs sont ceux du rapport, lus dans son fichier unique. Les DEUX noms —
% l'auteur et l'encadrant de stage — n'y sont pas écrits : ils viennent de
% `../report/noms.tex`, que le fichier des marqueurs charge s'il existe et qui
% n'est jamais commis. Sans lui, la planche de titre compose sans noms.
%
% Aucun texte fixe n'entoure un marqueur qui peut être vide : la ligne de
% l'encadrant passe par `\siRenseigne` et disparaît avec sa valeur, comme sur la
% page de garde du rapport.
%
% AUCUN NOMBRE N'EST TAPÉ ICI. Chaque valeur vient du registre des chiffres, par son
% identifiant, exactement comme dans le rapport : `../report/chiffres.tex` est le
% fichier que le registre produit, et il porte `\chiffre` comme les valeurs. Sans
% cette règle, la présentation serait le seul document du projet où un chiffre est
% écrit à la main — et c'est par cette voie que cinq valeurs du projet sont devenues
% fausses. Le contrôle du registre a vu sa portée étendue à ce répertoire pour que
% la règle soit tenue et non seulement écrite.
%
% Ce chapitre repose sur :
%   CHI: ablation-comparaisons-variante-a, ablation-ecart-complet, ablation-ecart-variante-a, ablation-ecart-variante-c, ablation-f-variante-c, admissions-2024
%   CHI: baseline-f-mesure, capacite-fonctionnelle, collision-groupes-nom, collision-groupes-piece, collision-vides-piece, colonnes-source-obs
%   CHI: colonnes-source-registre, comparaisons-modele, consultations-2024, defaut-taux-doublons, dim-patient-fiches-plusieurs-versions, dim-patient-lignes
%   CHI: dimensions, examens-2024, facteur-annualisation, faits, grappes-exactes-global, grappes-exactes-restreint
%   CHI: identifiants-distincts, interventions-2024, journees-2024, paires-candidates, paires-possibles, parametres-obs
%   CHI: parametres-total, periode-jours, personnes-estimees, relation-brute-ensemble, relation-inter-ensemble, relation-intra-ensemble
%   CHI: relations-injectees, relations-non-reprises, rotation-affichee, rotation-affichee-avant, seuil-retenu, source-lignes-total
%   CHI: source-tables, sur-fusionnes-global, taches-graphe, tdb-graphiques, tdb-indicateurs, tdb-pages
%   CHI: tdb-pages-methode, tdb-pages-pilotage, trot-publie, vt-paires-injectees
%   SER: (aucune)

\documentclass[aspectratio=169]{beamer}

\usepackage[T1]{fontenc}
\usepackage[utf8]{inputenc}
\usepackage{lmodern}
\usepackage[french]{babel}
\usepackage{csquotes}
\usepackage{xcolor}
\usepackage{tikz}

\usetheme{default}
\usecolortheme{default}
\setbeamertemplate{navigation symbols}{}
\setbeamertemplate{footline}[frame number]

% Les styles de tracé de la présentation.
%
% Ils sont écrits ici plutôt qu'empruntés à `report/images.tex` : ce fichier-là porte l'apparat des
% figures du rapport — emplacements réservés, logotypes, compteur de captures — dont une planche
% n'a que faire, et il déclare ses couleurs par des noms que le document principal du rapport
% définit. Une planche qui l'inclurait chargerait un apparat inutile et dépendrait d'un fichier
% qu'elle ne compose pas.
%
% La couleur d'accent est celle de la page de garde du rapport, à l'identique : le support de
% soutenance et le document remis se ressemblent, sans que l'un lise l'autre.

\definecolor{accentgarde}{RGB}{23,54,93}
\definecolor{griscartouche}{gray}{0.42}

\setbeamercolor{structure}{fg=accentgarde}
\setbeamercolor{frametitle}{fg=accentgarde}
\setbeamercolor{title}{fg=accentgarde}
\setbeamerfont{frametitle}{series=\bfseries}
\setbeamertemplate{itemize item}{\color{accentgarde}\textbullet}
\setbeamertemplate{itemize subitem}{\color{griscartouche}--}

\usetikzlibrary{positioning, shapes.geometric, fit, calc, arrows.meta}

\tikzset{
  % — une brique de couche, ou une étape.
  briquePres/.style={
    rectangle, draw=accentgarde, line width=0.6pt, align=center,
    font=\scriptsize, inner sep=2.5pt, minimum height=7mm, text width=24mm,
  },
  briqueLargePres/.style={briquePres, text width=54mm},
  % — un cadre secondaire : trait pointillé.
  briqueOptPres/.style={briquePres, dash pattern=on 1.6pt off 1.6pt},
  % — l'intitulé d'une couche.
  couchePres/.style={font=\scriptsize\scshape, text=griscartouche, align=right},
  % — une association vers une table descriptive.
  assocPres/.style={draw=accentgarde, line width=0.5pt},
  % — une flèche de flux.
  fluxPres/.style={draw=accentgarde, line width=0.6pt, -{Latex[length=1.8mm,width=1.4mm]}},
  % — une accolade de regroupement.
  notePres/.style={font=\scriptsize\itshape, text=griscartouche},
}

% Les marqueurs nominatifs du rapport, et le seul endroit du dépôt où un nom de
% personne peut figurer. Tout autre fichier les emploie par leur commande, jamais
% en écrivant la valeur.
%
% Un marqueur laissé vide fait échouer un contrôle : c'est la propriété que le
% critère de terminaison du projet exige, et tests/test_marqueurs_nominatifs.py
% la vérifie sur CE fichier, jamais sur le PDF produit.
%
% Renseigner une valeur, c'est écrire entre les accolades de la commande. Ne pas
% renommer les commandes : le contrôle les cherche par leur nom.

% --- l'état du document -------------------------------------------------------
% Deux valeurs, et deux seulement.
%
%   brouillon : le rapport n'est pas remis. Les cinq marqueurs nominatifs
%               ci-dessous sont TOUS vides, et la page de garde porte une mention
%               visible de version de travail.
%   remise    : le rapport est remis. Les cinq marqueurs sont TOUS renseignés, et
%               la mention disparaît.
%
% Un contrôle vérifie l'accord entre cet état et les cinq marqueurs : renseigner
% quatre marqueurs sur cinq est rouge dans les deux états, et c'est l'oubli qui
% arrive. Passer à la remise ne coûte que ce mot — et le contrôle dira alors quel
% marqueur manque encore.
\newcommand{\etatDuDocument}{brouillon}

% --- marqueurs nominatifs : ils portent un nom de personne ---------------------
\newcommand{\marqueurAuteur}{}
\newcommand{\marqueurEncadrantAcademique}{}
\newcommand{\marqueurEncadrantProfessionnel}{}
\newcommand{\marqueurPresidentJury}{}
\newcommand{\marqueurExaminateur}{}

% --- marqueurs de contexte : ils ne portent aucun nom de personne --------------
% Le contrôle ne les exige pas renseignés ; ils sont ici pour que la page de
% garde n'écrive aucune valeur en dur.
\newcommand{\marqueurEtablissementFormation}{}
\newcommand{\marqueurFiliere}{}
\newcommand{\marqueurAnneeUniversitaire}{}
\newcommand{\marqueurOrganismeAccueil}{Hôpital Sidi Saïd, Meknès}
\newcommand{\marqueurDateSoutenance}{}

% --- ce que l'état fait porter au document ------------------------------------
% Le mécanisme est typographique : le document se suffit, et ne dépend d'aucun
% contrôle pour dire ce qu'il est. `\siBrouillon{...}` compose son argument en
% état de brouillon, et rien en état de remise.
%
% La comparaison passe par \ifx sur deux commandes développées : c'est la forme
% que le noyau offre sans paquet supplémentaire, et elle compare les définitions
% caractère pour caractère.
\newcommand{\etatBrouillonAttendu}{brouillon}
\newcommand{\siBrouillon}[1]{%
  \ifx\etatDuDocument\etatBrouillonAttendu#1\fi%
}

% La mention elle-même, écrite une fois et employée par les deux documents.
\newcommand{\mentionVersionDeTravail}{%
  \siBrouillon{\textbf{Version de travail} — document non remis}%
}

% Fichier produit mécaniquement à partir de docs/chiffres/registre_chiffres.yml.
% Ne pas modifier à la main : docs/chiffres/generer_chiffres_tex.py le réécrit.
%
% Le rapport n'écrit aucune valeur en clair. Il écrit \chiffre{identifiant}, et la valeur vient
% d'ici, mesurée par la commande que le registre porte en regard de cet identifiant.
%
% Un identifiant inconnu déclenche une erreur de compilation qui le nomme : un chiffre absent doit
% arrêter la composition, jamais laisser un blanc dans une phrase.

\newcommand{\chiffre}[1]{%
  \ifcsname chiffre@#1\endcsname
    \csname chiffre@#1\endcsname
  \else
    \GenericError{}{Chiffre inconnu : #1}{}{Cet identifiant n'existe pas au registre des chiffres.}%
  \fi
}


\expandafter\def\csname chiffre@source-patients-lignes\endcsname{29\,107}
\expandafter\def\csname chiffre@source-rendez-vous-lignes\endcsname{14\,169}
\expandafter\def\csname chiffre@source-passages-lignes\endcsname{40\,650}
\expandafter\def\csname chiffre@source-passages-urgences-lignes\endcsname{27\,360}
\expandafter\def\csname chiffre@source-mouvements-lignes\endcsname{6\,310}
\expandafter\def\csname chiffre@source-prises-en-charge-lignes\endcsname{16\,430}
\expandafter\def\csname chiffre@source-factures-lignes\endcsname{21\,066}
\expandafter\def\csname chiffre@source-lignes-facture-lignes\endcsname{160\,936}
\expandafter\def\csname chiffre@source-encaissements-lignes\endcsname{23\,231}
\expandafter\def\csname chiffre@source-creances-lignes\endcsname{5\,876}
\expandafter\def\csname chiffre@source-relances-lignes\endcsname{1\,014}
\expandafter\def\csname chiffre@source-lignes-total\endcsname{346\,149}
\expandafter\def\csname chiffre@source-tables\endcsname{11}
\expandafter\def\csname chiffre@quarantaine-patients-lignes\endcsname{64}
\expandafter\def\csname chiffre@quarantaine-rendez-vous-lignes\endcsname{42}
\expandafter\def\csname chiffre@quarantaine-lignes-total\endcsname{106}
\expandafter\def\csname chiffre@colonnes-source-registre\endcsname{175}
\expandafter\def\csname chiffre@colonnes-source-obs\endcsname{81}
\expandafter\def\csname chiffre@colonnes-source-doc\endcsname{72}
\expandafter\def\csname chiffre@colonnes-source-hyp\endcsname{22}
\expandafter\def\csname chiffre@colonnes-source-catalogue\endcsname{175}
\expandafter\def\csname chiffre@colonnes-source-catalogue-obs\endcsname{81}
\expandafter\def\csname chiffre@colonnes-source-catalogue-doc\endcsname{72}
\expandafter\def\csname chiffre@colonnes-source-catalogue-hyp\endcsname{22}
\expandafter\def\csname chiffre@colonnes-aval-total\endcsname{398}
\expandafter\def\csname chiffre@colonnes-aval-decrites\endcsname{398}
\expandafter\def\csname chiffre@colonnes-intermediate\endcsname{175}
\expandafter\def\csname chiffre@colonnes-marts\endcsname{223}
\expandafter\def\csname chiffre@parametres-total\endcsname{256}
\expandafter\def\csname chiffre@parametres-fichiers\endcsname{25}
\expandafter\def\csname chiffre@parametres-obs\endcsname{5}
\expandafter\def\csname chiffre@parametres-doc\endcsname{52}
\expandafter\def\csname chiffre@parametres-hyp\endcsname{199}
\expandafter\def\csname chiffre@capacite-fonctionnelle\endcsname{40}
\expandafter\def\csname chiffre@journees-2024\endcsname{7\,851}
\expandafter\def\csname chiffre@admissions-2024\endcsname{1\,197}
\expandafter\def\csname chiffre@consultations-2024\endcsname{4\,142}
\expandafter\def\csname chiffre@prelevements-2024\endcsname{9\,625}
\expandafter\def\csname chiffre@examens-2024\endcsname{49\,597}
\expandafter\def\csname chiffre@interventions-2024\endcsname{0}
\expandafter\def\csname chiffre@tom-publie\endcsname{53,8}
\expandafter\def\csname chiffre@dms-publie\endcsname{6,6}
\expandafter\def\csname chiffre@periode-jours\endcsname{912}
\expandafter\def\csname chiffre@periode-debut\endcsname{2024-01-01}
\expandafter\def\csname chiffre@periode-fin\endcsname{2026-06-30}
\expandafter\def\csname chiffre@sejours-produits\endcsname{2\,980}
\expandafter\def\csname chiffre@consultations-produites\endcsname{10\,310}
\expandafter\def\csname chiffre@passages-urgences-produits\endcsname{27\,360}
\expandafter\def\csname chiffre@patients-produits\endcsname{29\,107}
\expandafter\def\csname chiffre@personnes-distinctes\endcsname{25\,842}
\expandafter\def\csname chiffre@urgences-scenario-retenu\endcsname{30}
\expandafter\def\csname chiffre@urgences-borne-basse\endcsname{14}
\expandafter\def\csname chiffre@urgences-scenario-median\endcsname{30}
\expandafter\def\csname chiffre@urgences-borne-haute\endcsname{54}
\expandafter\def\csname chiffre@defaut-taux-doublons\endcsname{0,04}
\expandafter\def\csname chiffre@defaut-dates-aberrantes\endcsname{0,003}
\expandafter\def\csname chiffre@defaut-ages-incoherents\endcsname{0,005}
\expandafter\def\csname chiffre@defaut-rdv-doublon-creneau\endcsname{0,012}
\expandafter\def\csname chiffre@defaut-factures-sans-pec\endcsname{0,03}
\expandafter\def\csname chiffre@defaut-familles\endcsname{13}
\expandafter\def\csname chiffre@part-sexe-masculin\endcsname{0,497}
\expandafter\def\csname chiffre@part-urbain\endcsname{0,639}
\expandafter\def\csname chiffre@part-sans-couverture\endcsname{0,2}
\expandafter\def\csname chiffre@part-couverture-principale\endcsname{0,5}
\expandafter\def\csname chiffre@tranches-age\endcsname{6}
\expandafter\def\csname chiffre@niveaux-tri\endcsname{5}
\expandafter\def\csname chiffre@part-tri-niveau-4\endcsname{0,41}
\expandafter\def\csname chiffre@part-tri-niveau-5\endcsname{0,22}
\expandafter\def\csname chiffre@part-sejours-depuis-urgences\endcsname{0,57}
\expandafter\def\csname chiffre@mouvements-lits-distincts\endcsname{40}
\expandafter\def\csname chiffre@factures-avec-creance\endcsname{5\,486}


\title{Chaîne de données et tableau de bord\\pour un établissement hospitalier}
\author{\marqueurAuteur}
% Le seul texte fixe de la planche de titre est le saut de ligne qui sépare
% l'établissement de la filière, et il ne compose que si les DEUX portent une
% valeur : un séparateur seul en tête ou en queue serait une ligne blanche
% inexpliquée. Les trois appels disent exactement cela, sans commande de plus.
\institute{%
  \siRenseigne{\marqueurEtablissementFormation}{\marqueurEtablissementFormation}%
  \siRenseigne{\marqueurEtablissementFormation}{\siRenseigne{\marqueurFiliere}{\\}}%
  \siRenseigne{\marqueurFiliere}{\marqueurFiliere}%
}
\date{\marqueurDateSoutenance}

\begin{document}

% ---------------------------------------------------------------- 1. Titre
\begin{frame}[plain]
  \siBrouillon{\begin{center}\fbox{\mentionVersionDeTravail}\end{center}}
  \titlepage
  \begin{center}
    \small
    \begin{tabular}{rl}
      \siRenseigne{\marqueurEncadrantProfessionnel}{Encadrant de stage & \marqueurEncadrantProfessionnel \\}
      Organisme d'accueil & \marqueurOrganismeAccueil \\
    \end{tabular}
  \end{center}
\end{frame}

\section{Contexte et problème}

% ---------------------------------------------------------------- 2. Le service
\begin{frame}{Un profil du logiciel sert une mission que le règlement place ailleurs}
  \begin{columns}[T]
    \begin{column}{0.52\textwidth}
      \small
      Le règlement intérieur des hôpitaux prescrit \textbf{neuf} missions au service d'accueil
      et d'admission. Le logiciel en propose \textbf{cinq} profils.
      \medskip

      Quatre missions reçoivent un profil. \textbf{Cinq n'en reçoivent aucun.}
      \medskip

      Et le profil de recouvrement ne répond à aucune mission de l'article~35 : le recouvrement
      relève de l'article~9, qui en charge un autre pôle.
    \end{column}
    \begin{column}{0.44\textwidth}
      \begin{tikzpicture}[x=1mm,y=1mm]
        \node[briquePres, text width=30mm] (a) at (0,0) {Neuf missions\\prescrites};
        \node[briquePres, text width=30mm] (b) at (0,-16) {Cinq profils\\applicatifs};
        \draw[fluxPres] (a) -- node[notePres,right,inner sep=1.5pt] {quatre liens} (b);
        \node[notePres, align=center, text width=34mm] at (0,-27)
          {cinq missions sans profil\\un profil sans mission};
      \end{tikzpicture}
    \end{column}
  \end{columns}
\end{frame}

% ---------------------------------------------------------------- 3. Au poste
\begin{frame}{Un champ renseigné n'est pas un champ juste}
  \small
  Au poste d'accueil, un agent saisit parfois un numéro de téléphone inventé, lorsque la personne
  en face ne connaît pas le sien.
  \medskip

  Il ne fabrique pas une erreur : \textbf{il fait passer un dossier avec ce qu'il a.}
  \medskip

  Conséquence pour toute la chaîne :
  \begin{itemize}
    \item le taux de renseignement d'une colonne mesure ce qui a été \emph{saisi} ;
    \item il ne mesure jamais ce qui est \emph{vrai} ;
    \item un rappel envoyé sur un tel numéro part, et personne ne le reçoit.
  \end{itemize}
\end{frame}

% ---------------------------------------------------------------- 4. Le problème
\begin{frame}{Le service saisit tout, et n'exploite presque rien}
  \small
  Le service enregistre chaque jour l'accueil, les rendez-vous, les passages, les mouvements et
  la facturation. Cette matière existe.
  \medskip

  \textbf{Elle ne sert presque à rien tant qu'aucun chemin ne la mène jusqu'à un indicateur.}
  \medskip

  Établir les statistiques et gérer l'information hospitalière est pourtant une mission
  \emph{propre} du service — la cinquième de l'article~35 —, et non une fonction déléguée à une
  cellule tierce.
\end{frame}

% ---------------------------------------------------------------- 5. La contrainte
\begin{frame}{Aucune donnée réelle ne sort du service}
  \small
  C'est la contrainte de cadrage, et elle est dite d'emblée parce qu'elle gouverne tout le reste.
  \medskip

  \begin{itemize}
    \item Aucun extrait, aucun échantillon, aucune donnée anonymisée.
    \item La chaîne est donc construite sur un \textbf{jeu de données engendré}.
    \item Elle lit des fichiers, non ce générateur : ce qui changerait en la branchant sur des
      données réelles est écrit au rapport, chapitre par chapitre.
  \end{itemize}
  \medskip

  \textbf{Ce que la chaîne démontre n'est donc jamais un constat sur le service.} La distinction
  revient à chaque planche qui porte un résultat.
\end{frame}

\section{Le jeu de données}

% ---------------------------------------------------------------- 6. Les volumes
\begin{frame}{L'établissement est nommé dans le recueil statistique du ministère}
  \small
  Les volumes ne sont pas inventés : ils sont relevés dans une série annuelle, pour cet
  établissement, sous son nom.
  \medskip

  \begin{center}
    \begin{tabular}{@{}lr@{}}
      \hline
      Capacité litière fonctionnelle & \chiffre{capacite-fonctionnelle} lits \\
      Journées d'hospitalisation & \chiffre{journees-2024} \\
      Admissions & \chiffre{admissions-2024} \\
      Consultations & \chiffre{consultations-2024} \\
      Examens & \chiffre{examens-2024} \\
      Interventions chirurgicales & \chiffre{interventions-2024} \\
      \hline
    \end{tabular}
  \end{center}
  \medskip

  Le dernier chiffre n'est pas une donnée manquante : \textbf{l'établissement n'exerce aucune
  activité chirurgicale}, et deux appuis indépendants l'établissent.
\end{frame}

% ---------------------------------------------------------------- 7. L'asymétrie
\begin{frame}{La structure est observée, les valeurs ne le sont pas}
  \small
  \begin{columns}[T]
    \begin{column}{0.48\textwidth}
      \textbf{Ce que l'écran a montré}
      \medskip

      \chiffre{colonnes-source-obs} colonnes sur \chiffre{colonnes-source-registre} viennent
      d'un relevé d'écran, chacune sous un identifiant que le lecteur retrouve en annexe.
    \end{column}
    \begin{column}{0.48\textwidth}
      \textbf{Ce que l'écran n'a pas montré}
      \medskip

      \chiffre{parametres-obs} paramètres de génération sur \chiffre{parametres-total} reposent
      sur une observation. Les autres viennent d'une source publiée ou d'une hypothèse posée.
    \end{column}
  \end{columns}
  \bigskip

  \textbf{L'asymétrie est le fait central de ce travail.} Le \emph{modèle} est observé ; les
  \emph{valeurs} qui le remplissent ne le sont pas, et aucune conclusion ne peut donc porter sur
  le service.
\end{frame}

\section{La chaîne de données}

% ---------------------------------------------------------------- 8. L'architecture (longue)
\begin{frame}{L'architecture : aucune couche ne modifie la précédente}
  \begin{center}
    \begin{tikzpicture}[x=0.95mm,y=0.95mm]
      \node[briquePres, text width=22mm] (zone) at (0,0) {Zone\\d'atterrissage};
      \node[briquePres, text width=22mm] (src)  at (28,0) {\texttt{source}};
      \node[briqueOptPres, text width=22mm] (qua) at (28,-16) {\texttt{quarantaine}};
      \node[briquePres, text width=22mm] (int)  at (56,0) {\texttt{intermediate}};
      \node[briquePres, text width=22mm] (mar)  at (84,0) {\texttt{marts}};
      \node[briqueOptPres, text width=22mm] (lin) at (84,-16) {\texttt{linkage}};
      \node[briquePres, text width=22mm] (ins)  at (110,0) {\texttt{instantane}};

      \draw[fluxPres] (zone) -- (src);
      \draw[fluxPres] (src) -- (qua);
      \draw[fluxPres] (src) -- (int);
      \draw[fluxPres] (int) -- (mar);
      \draw[fluxPres] (mar) -- (lin);
      \draw[fluxPres] (mar) -- (ins);

      \node[notePres, anchor=north, text width=26mm, align=center] at (0,-8)
        {fichiers partitionnés};
      \node[notePres, anchor=north, text width=26mm, align=center] at (28,-24)
        {lignes rejetées,\\avec leur motif};
      \node[notePres, anchor=north, text width=26mm, align=center] at (56,-8)
        {typage,\\normalisation};
      \node[notePres, anchor=north, text width=26mm, align=center] at (84,-24)
        {rapprochement\\d'identités};
      \node[notePres, anchor=north, text width=28mm, align=center] at (110,-8)
        {tables figées\\que le tableau de bord lit};
    \end{tikzpicture}
  \end{center}
  \small
  \begin{itemize}
    \item La couche source est une \textbf{copie fidèle} des fichiers, défauts compris ; corriger à
      l'entrée aurait effacé ce qu'il fallait mesurer.
    \item \chiffre{taches-graphe} tâches en chaîne, ordonnancées chaque jour, et les
      \textbf{dimensions se construisent avant les faits}.
    \item Le rafraîchissement de l'instantané n'est pas une reconstruction : les tables neuves sont
      bâties à côté, puis \textbf{tous les noms sont échangés en une seule transaction}.
  \end{itemize}
\end{frame}

% ---------------------------------------------------------------- 9. Étoile et historisation
\begin{frame}{Une fiche modifiée fait une version, pas une personne}
  \small
  \chiffre{dimensions} dimensions, \chiffre{faits} faits. La dimension patient est
  \emph{historisée} : elle garde chaque état successif d'une fiche.
  \medskip

  \begin{center}
    \begin{tikzpicture}[x=1mm,y=1mm]
      \node[briquePres, text width=32mm] (v) at (0,0) {\chiffre{dim-patient-lignes} versions de fiche};
      \node[briquePres, text width=32mm] (i) at (44,0) {\chiffre{identifiants-distincts} identifiants};
      \node[briquePres, text width=32mm] (p) at (88,0) {\chiffre{personnes-estimees} personnes};
      \draw[fluxPres] (v) -- (i);
      \draw[fluxPres] (i) -- (p);
      \node[notePres, anchor=north, text width=40mm, align=center] at (22,-7)
        {\chiffre{dim-patient-fiches-plusieurs-versions} en portent plus d'une};
      \node[notePres, anchor=north, text width=40mm, align=center] at (66,-7)
        {des doublons d'identité les réunissent};
    \end{tikzpicture}
  \end{center}
  \medskip

  \textbf{Trois grandeurs, non deux.} Les confondre est l'erreur la plus facile à commettre sur ce
  jeu de données, et elle a été commise trois fois dans les documents de travail de ce projet.
\end{frame}

\section{Qualité et rapprochement d'identités}

% ---------------------------------------------------------------- 10. Pourquoi les doublons
\begin{frame}{Les doublons sont établis par l'outil et par l'audit, jamais par une observation}
  \small
  \begin{itemize}
    \item \textbf{Par l'outil.} L'écran de recherche porte un onglet de recherche \emph{pondérée}
      sur l'index maître des patients, et une colonne de \emph{probabilité} en regard de chaque
      résultat. Une recherche qui rend des probabilités n'est pas une recherche exacte.
    \item \textbf{Par l'audit.} Un rapport de contrôle consacré à l'établissement documente le
      phénomène.
  \end{itemize}
  \medskip

  \textbf{Et une réserve qui compte autant que le reste :} aucune recherche conduite au poste n'a
  rendu plusieurs résultats proches d'un même patient. \emph{Le problème n'a pas été constaté sur
  place.}
  \medskip

  Le jeu de données en injecte donc une part connue — \chiffre{defaut-taux-doublons} — et
  \chiffre{vt-paires-injectees} paires forment la vérité terrain.
\end{frame}

% ---------------------------------------------------------------- 11. Comment on les retrouve
\begin{frame}{Le blocage écarte, le modèle tranche}
  \small
  \begin{center}
    \begin{tikzpicture}[x=1mm,y=1mm]
      \node[briquePres, text width=34mm] (t) at (0,0) {\chiffre{paires-possibles} paires possibles};
      \node[briquePres, text width=34mm] (c) at (48,0) {\chiffre{paires-candidates} paires candidates};
      \node[briquePres, text width=34mm] (m) at (96,0) {\chiffre{comparaisons-modele} comparaisons};
      \draw[fluxPres] (t) -- node[notePres,above,inner sep=1.5pt] {blocage} (c);
      \draw[fluxPres] (c) -- node[notePres,above,inner sep=1.5pt] {modèle} (m);
    \end{tikzpicture}
  \end{center}
  \medskip

  \begin{itemize}
    \item \textbf{Le blocage} réunit quatre règles — nom et adresse, nom et téléphone, pièce
      d'identité, parents et naissance — et son rappel sur la vérité terrain est total.
    \item \textbf{Le modèle} est probabiliste : il estime, pour chaque comparaison, la probabilité
      d'accord chez les vrais couples et chez les faux.
    \item \textbf{Le seuil retenu} est \chiffre{seuil-retenu}.
  \end{itemize}
\end{frame}

% ---------------------------------------------------------------- 12. Ventilation (longue)
\begin{frame}{Ce que la ligne de base retrouve, et ce qu'elle manque}
  \small
  La ligne de base est la collision exacte : deux fiches sont dites identiques si un champ le dit
  à la lettre. Sa F-mesure vaut \chiffre{baseline-f-mesure}.
  \medskip

  \begin{center}
    \begin{tabular}{@{}lrl@{}}
      \hline
      Groupes par nom exact & \chiffre{collision-groupes-nom} & retrouvés sans modèle \\
      Groupes par pièce d'identité & \chiffre{collision-groupes-piece} & idem \\
      Fiches sans pièce d'identité & \chiffre{collision-vides-piece} & \textbf{invisibles à la collision} \\
      \hline
      Grappes sur la population restreinte & \chiffre{grappes-exactes-restreint} & \\
      Grappes sur la population entière & \chiffre{grappes-exactes-global} & \\
      Fiches sur-fusionnées & \chiffre{sur-fusionnes-global} & \\
      \hline
    \end{tabular}
  \end{center}
  \medskip

  \textbf{Le chiffre qui compte est le troisième.} Une fiche dont la pièce d'identité est vide ne
  peut pas entrer dans une collision exacte : c'est exactement la population que le modèle
  probabiliste existe pour retrouver, et c'est aussi celle que le service saisit le moins bien.
\end{frame}

% ---------------------------------------------------------------- 13. Ablation (longue)
\begin{frame}{L'ablation : la F-mesure ne bouge pas, la marge s'effondre}
  \small
  Retirer des comparaisons au modèle et regarder la F-mesure ne dit rien. Il faut regarder
  \emph{l'écart de poids} entre les paires appariées et les autres.
  \medskip

  \begin{center}
    \begin{tabular}{@{}lrr@{}}
      \hline
      Variante & Comparaisons & Écart de poids \\
      \hline
      Modèle complet & \chiffre{comparaisons-modele} & \chiffre{ablation-ecart-complet} \\
      Variante A & \chiffre{ablation-comparaisons-variante-a} & \chiffre{ablation-ecart-variante-a} \\
      Variante C & \chiffre{comparaisons-modele} & \chiffre{ablation-ecart-variante-c} \\
      \hline
    \end{tabular}
  \end{center}
  \medskip

  La variante C garde une F-mesure de \chiffre{ablation-f-variante-c} — quasi parfaite — et son
  écart de poids est pourtant \textbf{négatif}.
  \medskip

  \textbf{Un modèle qui classe encore juste peut avoir cessé de discriminer.} La F-mesure seule
  l'aurait laissé passer.
\end{frame}

% ---------------------------------------------------------------- 14. Limites de l'évaluation
\begin{frame}{Ce que cette évaluation ne prouve pas}
  \small
  \begin{itemize}
    \item \textbf{La vérité terrain est connue parce qu'elle a été injectée.} Sur des données
      réelles, cette connaissance n'existe pas : il faudrait un échantillon annoté à la main, avec
      son coût et son incertitude propres.
    \item \textbf{Précision et rappel cessent alors d'être mesurables tels quels.}
    \item \textbf{Les défauts sont ceux que j'ai semés}, à un taux que j'ai choisi. Un fichier réel
      en porte d'autres, dans d'autres proportions.
  \end{itemize}
  \medskip

  Ce que l'évaluation établit est plus modeste et plus solide : \textbf{la chaîne sait produire
  ces grandeurs, à la maille où une décision se prend.}
\end{frame}

\section{Analyse de l'activité}

% ---------------------------------------------------------------- 15. Paradoxe de composition
\begin{frame}{Trois corrélations, et une seule d'entre elles ment}
  \small
  Le lien entre le délai d'obtention d'un rendez-vous et l'absentéisme se mesure de trois façons,
  et les trois ne disent pas la même chose.
  \medskip

  \begin{center}
    \begin{tabular}{@{}lr@{}}
      \hline
      Entre activités & \chiffre{relation-inter-ensemble} \\
      À l'intérieur de chaque activité & \chiffre{relation-intra-ensemble} \\
      Sur les rendez-vous pris ensemble & \chiffre{relation-brute-ensemble} \\
      \hline
    \end{tabular}
  \end{center}
  \medskip

  \textbf{La première est négative, la deuxième positive, la troisième presque nulle.} Publier
  l'une des trois seule serait un artefact de composition : le tableau de bord porte les deux
  premières côte à côte, et jamais l'une sans l'autre.
\end{frame}

\section{Le tableau de bord}

% ---------------------------------------------------------------- 16. Deux sections
\begin{frame}{Deux sections, et la seconde ne sert aucune décision}
  \small
  \chiffre{tdb-pages} pages, \chiffre{tdb-indicateurs} indicateurs, \chiffre{tdb-graphiques}
  tracés, sur \chiffre{periode-jours} jours et \chiffre{source-lignes-total} lignes réparties sur
  \chiffre{source-tables} tables.
  \medskip

  \begin{columns}[T]
    \begin{column}{0.48\textwidth}
      \textbf{Pilotage du service} — \chiffre{tdb-pages-pilotage} pages
      \medskip

      Chaque page sert une décision : ouvrir des créneaux, dimensionner le tri, répartir les lits,
      prioriser le recouvrement.
    \end{column}
    \begin{column}{0.48\textwidth}
      \textbf{Évaluation de la chaîne} — \chiffre{tdb-pages-methode} pages
      \medskip

      Aucune décision de service ne s'y adosse. Ces pages décrivent un modèle et une provenance,
      non une file d'attente.
    \end{column}
  \end{columns}
  \medskip

  \textbf{Aucun indicateur circulaire n'est classé en pilotage}, et un contrôle le vérifie.
\end{frame}

% ---------------------------------------------------------------- 17. L'indicateur faux
\begin{frame}{Un indicateur réglementaire faux, qu'aucun contrôle ne pouvait voir}
  \small
  Le taux de rotation affichait \chiffre{rotation-affichee-avant} là où la statistique nationale
  publie \chiffre{trot-publie} : faux d'un facteur \chiffre{facteur-annualisation}, soit exactement
  le rapport de la durée de la période à l'année de référence. \textbf{La page n'annualisait pas.}
  \medskip

  \textbf{Le contrôle qui le surveillait était vert, et il l'est resté pendant toute la durée du
  défaut.} Il retranscrivait la formule de la page et se comparait à elle : une comparaison vraie
  par construction, incapable de départager une formule juste d'une formule fausse.
  \medskip

  Le contrôle qui l'a remplacé ne retranscrit rien : il appelle la page et confronte ses sorties à
  deux références qui ne peuvent pas bouger avec elle. Le taux vaut désormais
  \chiffre{rotation-affichee}.
  \medskip

  \textbf{Une seconde mesure n'a de valeur que si elle peut être fausse quand la première l'est.}
\end{frame}

\section{Limites et perspectives}

% ---------------------------------------------------------------- 18. Recommandations
\begin{frame}{Trois actions, et ce que chacune suppose}
  \small
  \begin{itemize}
    \item \textbf{Rappeler les patients} dont le rendez-vous approche, en commençant par les
      activités où le délai est le plus long. \emph{Suppose un numéro juste — et c'est là que la
      difficulté se trouve.}
    \item \textbf{Ouvrir une filière courte} pour les deux niveaux de tri les moins graves.
      \emph{Suppose un tri formalisé dont le résultat est enregistré.}
    \item \textbf{Prioriser le recouvrement} par ancienneté et par débiteur. \emph{Suppose que la
      créance porte son débiteur, ce que le modèle observe.}
  \end{itemize}
  \medskip

  Aucun effort n'est chiffré et aucun effet n'est quantifié. Ce qui établit le constat vient
  toujours d'une source extérieure au jeu de données ; ce que l'indicateur démontre est seulement
  que la chaîne sait produire le nombre.
\end{frame}

% ---------------------------------------------------------------- 19. Limites
\begin{frame}{Les limites, énoncées et non concédées}
  \small
  \begin{itemize}
    \item \textbf{Le jeu de données n'est pas la réalité.} \chiffre{relations-injectees} relations
      y sont injectées, dont \chiffre{relations-non-reprises} qu'aucune conclusion ne reprend, et
      chaque conclusion dit laquelle elle reproduit.
    \item \textbf{L'observation a porté sur un poste et sur un profil}, une journée. Les quatre
      autres profils exigent des habilitations qu'un stagiaire ne reçoit pas.
    \item \textbf{La composition du service est inconnue} : aucune source consultée ne la donne, et
      ce rapport ne la suppose pas.
    \item \textbf{Aucun argument de performance} n'est tiré de la volumétrie : elle est celle d'un
      établissement de cette taille, non celle d'une épreuve de charge.
  \end{itemize}
\end{frame}

% ---------------------------------------------------------------- 20. Ensuite
\begin{frame}{Ce que je ferais ensuite}
  \small
  \begin{itemize}
    \item \textbf{Brancher la chaîne sur une extraction réelle}, sous convention d'accès — ce qui
      change de nature est écrit : la vérité terrain disparaît, les distributions cessent d'être
      posées, et un résultat cesse de démontrer pour se mettre à établir.
    \item \textbf{Mesurer le taux d'aboutissement des rappels} avant d'en attendre un effet : c'est
      lui qui dira quelle part du fichier de contact est utilisable.
    \item \textbf{Écrire au fil de l'eau plutôt qu'à la fin.} J'ai découvert en rédigeant des
      erreurs qu'aucun contrôle n'avait vues — écrire est aussi une manière de vérifier.
  \end{itemize}
\end{frame}

% ---------------------------------------------------------------- 21. Merci
\begin{frame}[plain]
  \begin{center}
    \vspace{1.2cm}
    {\Large\bfseries\color{accentgarde} Merci de votre attention}
    \vspace{0.8cm}

    \small
    \begin{tabular}{rl}
      \siRenseigne{\marqueurEncadrantProfessionnel}{Encadrant de stage & \marqueurEncadrantProfessionnel \\}
      Organisme d'accueil & \marqueurOrganismeAccueil \\
    \end{tabular}
  \end{center}
\end{frame}

\end{document}
